\documentclass[11pt]{article}
\usepackage[margin=1in]{geometry}
\usepackage{amsmath}
\usepackage[T1]{fontenc}
\usepackage{lmodern}
\usepackage{enumitem}
\usepackage{microtype}

\title{Homework Answers}
\author{}
\date{}

\begin{document}
\maketitle

\section*{1. Bernie vs. a 150 Mbps transmission line}

Using \(1\,\mathrm{GB}=10^9\) bytes, Bernie carries
\[
3 \times 7\,\mathrm{GB} \times 8 = 168\,\mathrm{Gb}.
\]

At \(18\,\mathrm{km/h}=0.005\,\mathrm{km/s}\), the travel time for distance \(d\) is \(d/0.005\) s. Bernie is faster when
\[
\frac{168\times10^9}{d/0.005} > 150\times10^6
\quad\Rightarrow\quad d<5.6\,\mathrm{km}.
\]

Thus, Bernie has the higher data rate for \(0<d<5.6\,\mathrm{km}\) (one-way distance).

\begin{itemize}[leftmargin=*]
  \item (i) If Bernie's speed doubles: \(0<d<11.2\,\mathrm{km}\).
  \item (ii) If each tape's capacity doubles: \(0<d<11.2\,\mathrm{km}\).
  \item (iii) If the line rate doubles to 300 Mbps: \(0<d<2.8\,\mathrm{km}\).
\end{itemize}

\section*{2. Layered protocols}

Two reasons:
\begin{enumerate}[leftmargin=*]
  \item They divide a complex networking task into manageable modules.
  \item A layer can be changed or replaced without redesigning the other layers.
\end{enumerate}

Disadvantage: layering can add overhead and reduce performance.

\section*{3. Retransmissions}

Each attempt uses one transmission; with probability \(p\), another attempt is needed. Assuming independent errors and \(p<1\),
\[
E[N]=1+pE[N]
\quad\Rightarrow\quad
E[N]=\frac{1}{1-p}.
\]
The mean number of transmissions is \(\frac{1}{1-p}\).

\section*{4. TCP vs. UDP}

\textbf{TCP} is connection-oriented and provides reliable, ordered delivery. \textbf{UDP} is connectionless and provides no delivery, ordering, or retransmission guarantees.

\section*{5. Length of one bit in 802.3}

The propagation speed is
\[
v=\frac{2}{3}(3\times10^8)=2\times10^8\,\mathrm{m/s}.
\]
At 10 Mbps, one bit lasts \(1/(10\times10^6)\) s, so
\[
L=v\cdot\frac{1}{10\times10^6}=20\,\mathrm{m}.
\]
One bit was \(20\) meters long.

\section*{6. International standards}

\begin{itemize}[leftmargin=*]
  \item Advantage: they ensure interoperability between equipment from different countries and vendors.
  \item Disadvantage: reaching agreement can be slow, and the resulting standards may be difficult to update.
\end{itemize}

\section*{7. Common OSI and TCP/IP layers}

Both models have \textbf{application}, \textbf{transport}, and \textbf{network} layers (called \textbf{Internet} in TCP/IP). In the four-layer TCP/IP model, the \textbf{network access} layer covers OSI's data link and physical layers.

\end{document}
